\begin{solapartat}
\punts{0,25} L'energia del fotó incident és:
\[ E_{fotons} = hf = h\,\frac{c}{\lambda} = 6,63\cdot 10^{-34}\,\frac{3,00\times 10^{8}}{460\times 10^{-9}} = 4,32\times 10^{-19}\ \mathrm{J} = 2,70\ \mathrm{eV} \]
\punts{0,5} L'energia cinètica màxima que poden adquirir els electrons emesos per efecte fotoelèctric depèn de l'energia dels fotons $E_{fotons}$ i la funció de treball del metall, $W$, de la manera següent:
\[ E_{C,\text{màx}} = E_{fotons} - W \]
El potencial de frenada és el potencial elèctric mínim que cal aplicar entre el metall emissor i l'ànode receptor per aturar completament els electrons emesos. Això impedeix que els electrons arribin a l'ànode. Per tant, el potencial de frenada, $V_f$, es relaciona amb l'energia cinètica màxima dels electrons:
\[ E_{C,\text{màx}} = eV_f \]
Combinant les dues expressions anteriors tenim:
\[ eV_f = E_{fotons} - W \Longrightarrow W = E_{fotons} - eV_f \]
I si substituïm:
\begin{align*}
W = E_{fotons} - eV_f &= 4,32\times 10^{-19} - 0,5\times 1,602\times 10^{-19}\\
&= 3,52\times 10^{-19}\ \mathrm{J} = 2,20\ \mathrm{eV}
\end{align*}
No és necessari que l'estudiant doni el resultat en J i eV. També es poden fer els càlculs directament en eV.

\medskip
\punts{0,5} L'energia cinètica màxima és:
\[ E_{C,\text{màx}} = E_{fotons} - W = 8,01\cdot 10^{-20}\ \mathrm{J} \]
I la velocitat d'aquests electrons és:
\[ v = \sqrt{\frac{2\,E_{C,\text{màx}}}{m_e}} = \sqrt{\frac{2\times 8,01\times 10^{-20}}{9,11\times 10^{-31}}} = 4,19\times 10^{5}\ \mathrm{m/s} \]
\end{solapartat}

\begin{solapartat}
\punts{0,25} A partir de les expressions anteriors, podem calcular el potencial de frenada com:
\[ V_f = \frac{E_{fotons} - W}{e} \]
\punts{0,5} Per a la llum verda:
\[ E_{fotons} = hf = h\,\frac{c}{\lambda} = 6,63\cdot 10^{-34}\,\frac{3,00\times 10^{8}}{520\times 10^{-9}} = 3,82\cdot 10^{-19}\ \mathrm{J} = 2,39\ \mathrm{eV} \]
\[ V_f = \frac{E_{fotons} - W}{e} = \frac{3,82\times 10^{-19} - 3,52\times 10^{-19}}{1,602\times 10^{-19}} = 0,19\ \mathrm{V} \]
\punts{0,5} Per a la llum vermella:
\[ E_{fotons} = hf = h\,\frac{c}{\lambda} = 6,63\cdot 10^{-34}\,\frac{3,00\times 10^{8}}{650\times 10^{-9}} = 3,06\cdot 10^{-19}\ \mathrm{J} = 1,91\ \mathrm{eV} \]
Tenim que $E_{fotons} < W$; els fotons no tenen prou energia per arrencar electrons del càtode. Per tant, no es produeix efecte fotoelèctric i no cal aplicar cap potencial de frenada.
\end{solapartat}
